%! AP Statistics — Unit 5, Lesson 5.3 — Lesson Plan
\documentclass[10pt]{article}
\usepackage{apstats-article}
\usepackage{apstats-boxes}

\begin{document}

\pageheader{Unit 5, Lesson 5.3}{Lesson Plan}

\begin{objectivebox}
\textbf{Topic:} 5.3 — The Central Limit Theorem\\
\textbf{Enduring Understanding:} UNC-3\\
\textbf{Learning Objective:} UNC-3.H — Describe the sampling distribution of a sample statistic.\\
\textbf{Essential Knowledge:} UNC-3.H.1–H.5\\
\textbf{Mathematical Practice:} 3.C — Verify that inference procedures apply in a given situation.
\end{objectivebox}

\begin{learningtargetbox}
By the end of this lesson, students will be able to:
\begin{itemize}
  \item State the Central Limit Theorem and identify its conditions.
  \item Describe the shape, center, and spread of a sampling distribution of $\bar{x}$.
  \item Determine when it is appropriate to use a normal model for a sampling distribution.
\end{itemize}
\end{learningtargetbox}

\begin{vocabbox}
Central Limit Theorem (CLT), sampling distribution of $\bar{x}$, standard error, independence condition, large counts condition ($n \ge 30$)
\end{vocabbox}

\begin{hookbox}
\textbf{Hook — Candy Machine Mystery (5 min):}
A candy machine fills bags with an average of 50 candies ($\mu = 50$) but with high variability ($\sigma = 12$). A quality-control inspector takes a sample of 36 bags each day. Why might the inspector's daily sample mean be much more predictable than any single bag count? Discuss with a partner.
\end{hookbox}

\begin{notesbox}
\textbf{Direct Instruction (15 min):}
\begin{enumerate}
  \item Review: sampling distribution concept from Lesson 5.1.
  \item Introduce CLT: shape of $\bar{x}$ distribution approaches normal as $n$ increases.
  \item State conditions: (1) random sample, (2) independence ($n \le 10\%$ of population), (3) $n \ge 30$ OR population is normal.
  \item Formulas: $\mu_{\bar{x}} = \mu$ and $\sigma_{\bar{x}} = \sigma/\sqrt{n}$.
  \item Connect to hook: show how $\sigma_{\bar{x}} = 12/\sqrt{36} = 2$ dramatically reduces variability.
\end{enumerate}
\end{notesbox}

\begin{practicebox}
\textbf{Guided Practice (15 min):} Students complete notes guided practice with teacher support; use cold-call questioning on condition checking.
\end{practicebox}

\begin{notesbox}
\textbf{Activity (15 min):} Tiered group activity — students apply CLT to real contexts, checking conditions and computing $\mu_{\bar{x}}$ and $\sigma_{\bar{x}}$.
\end{notesbox}

\begin{reflectionbox}
\textbf{Exit Ticket (5 min):} 3 quick items — identify conditions met/not met, compute standard error, interpret center.\\
\textbf{Homework:} Problems applying CLT; comparing individual vs.\ sampling distribution probabilities.\\
\textbf{AP Connection:} CLT is foundational for all inference in Units 6–9.
\end{reflectionbox}

\end{document}
